\chapter{Cipher Maps and Trapdoor Computing}
\label{ch:cipher-maps-trapdoor}

\begin{quote}
\textsc{After this chapter you will be able to:}
\begin{itemize}
  \item State the cipher-map abstraction: a function that computes
    faithfully on encrypted inputs up to bounded $\fprate$ and
    $\fnrate$.
  \item Recognize the Bernoulli framework as the error model for
    cipher maps; identify the chapter~4 Kronecker theorem as the
    composition law for chained cipher-map evaluations.
  \item Apply the chapter~7 PPV / NPV framework to confidentiality
    analysis: every cipher-map output has a posterior over plaintext
    truth, and that posterior is the confidentiality measure.
\end{itemize}
\end{quote}

% Section 17.1
% §17.1 What is a Cipher Map?
%
% Cross-reference resolution status:
%
% DEFINED here:
%   \label{sec:what-is-cipher-map}            -- this section
%   \label{def:cipher-map}                    -- the cipher-map definition
%
% ALREADY DEFINED (chapter 1):
%   \Cref{sec:noisy-bit}                      -- chapter 1 §1.2 fprate/fnrate
%
% ALREADY DEFINED (chapter 5):
%   \Cref{sec:bloom-from-scratch}             -- chapter 5 §5.2 Bloom filter
%
% FORWARD refs within chapter 17:
%   \Cref{sec:bernoulli-error-model}          -- §17.2 (Task 3)
%   \Cref{sec:fhe-sidebar}                    -- §17.4 (Task 5)
%
% §17.1 introduces the cipher-map abstraction: a function that computes
% faithfully on encrypted inputs up to bounded false-positive and
% false-negative rates. Five items: (1) plaintext setup; (2) cipher-map
% definition; (3) why approximate is acceptable; (4) the trapdoor key;
% (5) three example cipher maps.

\section{What is a Cipher Map?}
\label{sec:what-is-cipher-map}

The first sixteen chapters built a framework for noisy computation on
unencrypted data. The Bloom filter of chapter 5 returns wrong answers
because the structure is small; the classifier of chapter 9 returns
wrong answers because the world is hard; the sketch of chapter 10
returns wrong answers because the stream is large. In every case the
error is a property of the structure, and the user accepts it for
reasons of space or time.

This chapter takes the same Bernoulli machinery and applies it to a
different setting. The data is encrypted; the computation runs on the
encrypted form; only a holder of a secret key can read the result.
Bloom filters and noisy classifiers reappear, but now their false
positives and false negatives serve a second purpose alongside their
operational role: the bounded error is also what keeps the underlying
plaintext hidden from the party doing the computation.

\subsection*{The plaintext setting}

Begin with a function on plaintext data. Let $X$ be a domain of
plaintext values (records, sets, scalars, whatever the application
requires), let $Y$ be a codomain, and let
\[
  f : X \to Y
\]
be the function the user wants to compute. In the cleartext setting
the user holds an $x \in X$, applies $f$, and reads $f(x)$. Chapters
1 through 16 sometimes replaced $f$ with an approximate $\tilde f$
that occasionally returned the wrong $y$; the error analysis was
operational.

The cipher-map setting introduces a different actor. A holder of $x$
(the \emph{client}) wants $f(x)$ but does not want to give $x$ to the
machine doing the computation. Instead, the client encrypts $x$ under
an encryption map $\mathsf{Enc}$, sends the ciphertext to a
\emph{server}, and asks the server to compute something on
$\mathsf{Enc}(x)$ that the client can decrypt to obtain (something
close to) $f(x)$. The server never sees $x$. The server also never
sees $f(x)$ directly; what the server returns is itself an encryption
of an output, which the client decrypts using $\mathsf{Dec}$.

The shape of the computation is therefore not $f$ itself but a
companion map that operates on the encrypted representations. That
companion map is the cipher map.

\subsection*{The cipher-map abstraction}

\begin{definition}[Cipher map]
\label{def:cipher-map}
Let $f : X \to Y$ be a plaintext function. Let $\mathsf{Enc} : X \to
\mathsf{Enc}(X)$ be an encryption map onto an encrypted-domain space
$\mathsf{Enc}(X)$, with corresponding decryption $\mathsf{Dec} :
\mathsf{Enc}(Y) \to Y$. A \emph{cipher map} for $f$ is a function
\[
  \tilde f : \mathsf{Enc}(X) \to \mathsf{Enc}(Y)
\]
on the encrypted representations such that, for $x \in X$,
\[
  \Pr\left[\,\mathsf{Dec}(\tilde f(\mathsf{Enc}(x))) = f(x)\,\right]
  \;\ge\; 1 - \fnrate
\]
when $f(x)$ is a real positive output, and
\[
  \Pr\left[\,\mathsf{Dec}(\tilde f(\mathsf{Enc}(x))) = y \text{ for some
  spurious } y \ne f(x)\,\right] \;\le\; \fprate
\]
when $f(x)$ is the negative case. The pair $(\fprate, \fnrate)$ is
the cipher map's \emph{error profile}, in the Bernoulli sense of
chapter 1 \S\ref{sec:noisy-bit}.
\end{definition}

Three pieces of the definition deserve emphasis. The cipher map
operates entirely on the encrypted domain: every step from
$\mathsf{Enc}(x)$ through $\tilde f$ to $\mathsf{Enc}(y)$ leaves the
ciphertext encrypted. The decryption happens on the client side,
after the encrypted output has been returned. And the error rates
$\fprate, \fnrate$ are exactly the rates that the framework has
spent sixteen chapters analyzing: there is no new probabilistic
machinery to introduce; the cipher-map error is a Bernoulli error.

What is new is the second meaning carried by the error. A cipher map
with $\fprate = 0$ and $\fnrate = 0$ behaves identically to $f$ from
the client's view, but it also reveals more information about $x$ to
any adversary that observes $\tilde f$'s behavior across many
queries. A cipher map with positive error obscures $x$ in proportion
to the error magnitude. \Cref{sec:bernoulli-error-model} makes this
duality precise; \Cref{sec:trapdoor-paradigm} unpacks its
operational consequences.

\subsection*{Why approximate is acceptable}

It is reasonable to ask why one would accept bounded error in a
cryptographic setting. The cleartext approximate-set literature
tolerates Bloom-filter false positives because the alternative
(storing the full set with exact membership) costs too much memory;
the cipher-map literature tolerates cipher-map false positives for
an analogous reason in the cryptographic dimension.

The exact alternative exists. Fully homomorphic encryption (FHE),
introduced by Gentry in 2009 and refined by the BFV, BGV, and CKKS
schemes since, provides an $\tilde f$ with $\fprate = \fnrate = 0$
for any computable $f$. FHE is the cryptographer's answer to the
question ``can we compute exactly on encrypted data?'' and the
answer is yes. The catch is that FHE ciphertexts are large
(kilobytes to megabytes per encrypted bit in practical schemes) and
slow (seconds per encrypted bit operation). For workloads where the
plaintext computation runs in microseconds, an FHE realization can
run six orders of magnitude slower with three orders of magnitude
more memory.

Cipher maps trade error for practicality. A Bloom-filter-based
cipher map for set membership uses on the order of
$-\log_2 \fprate + \mu$ bits per inserted element (the optimal
formula of chapter 14), runs in constant time per query, and gives
a false-positive rate the client can tune. The trade-off is the
familiar Bernoulli one: space and runtime are exchanged for
controlled error. \Cref{sec:fhe-sidebar} returns to the
FHE-versus-cipher-map comparison in detail.

\subsection*{The trapdoor key}

The key that distinguishes a cipher map from any other noisy
function is the \emph{trapdoor key}: a secret that enables recovery
of $f(x)$ from $\tilde f(\mathsf{Enc}(x))$. The client holds the
trapdoor; the server does not. Without the trapdoor, the server
sees only opaque bit patterns flowing through $\tilde f$ and cannot
distinguish meaningful queries from filler, real outputs from noise,
or one plaintext value from another. With the trapdoor, the client
can decrypt the cipher-map output and recover the answer to within
the cipher map's $\fprate, \fnrate$ bounds.

Three properties make a trapdoor key worth its name. First, the
trapdoor is short: a key of constant length, not a representation
that grows with the data. Second, the trapdoor is necessary: no
polynomial-time computation without it can recover plaintext from
the cipher-map output. Third, the trapdoor is sufficient: a holder
of the key can decrypt every cipher-map output with the same per-call
cost as a plaintext function evaluation. These three properties
distinguish a cipher map from a generic encryption (which lacks
property three) and from a generic noisy function (which lacks
properties one and two).

\subsection*{Three example cipher maps}

Three cipher maps surface repeatedly in this chapter and in
chapter 18. Each is sketched here at the level needed to ground
\Cref{sec:bernoulli-error-model}'s analysis.

\textbf{Encrypted set membership.} Let $A \subset X$ be a plaintext
set the client wants to query without revealing to the server. The
function is the indicator $f_A(x) = \mathbf{1}[x \in A]$. The
client builds a Bloom filter on $A$ using a keyed hash family (the
key is the trapdoor), encrypts the resulting bit array under a
suitable scheme, and gives it to the server. To query, the client
encrypts the query element $x$, sends $\mathsf{Enc}(x)$ to the
server, and asks the server to test the encrypted bits. The server
returns a (possibly encrypted) bit. With probability at least
$1 - \fprate$ the bit is correct on negative queries
($x \notin A$); the bit is always correct on positive queries
($\fnrate = 0$), reflecting the one-sided error of the Bloom
construction.

\textbf{Encrypted equality test.} Two parties want to know whether
$x = y$ without revealing $x$ or $y$. The function is
$f_=(x, y) = \mathbf{1}[x = y]$. Under deterministic encryption,
identical plaintexts encrypt to identical ciphertexts, so equality
on ciphertexts coincides with equality on plaintexts and the cipher
map has $\fprate = \fnrate = 0$. The cost is a known information
leak: an adversary observing many equality tests can build the
distribution of plaintext values from ciphertext frequencies. This
is a cipher map with no Bernoulli error but with non-zero
information leak: a useful comparison case, since the framework's
$\fprate, \fnrate$ are not the only confidentiality knob.

\textbf{Encrypted ordering.} A client wants to know whether
$x \le y$ on encrypted inputs (for sorting, ranged queries, or
threshold tests). The function is $f_{\le}(x, y) =
\mathbf{1}[x \le y]$. Under order-preserving encryption, the
ciphertext encoding preserves the plaintext order, so the server
can compare ciphertexts directly. The cipher map has $\fprate =
\fnrate = 0$ on the comparison itself but leaks the entire
ranking of the underlying plaintexts to the server: an even larger
leak than the equality case. Order-preserving encryption is again
a cipher map with no Bernoulli error and a tightly characterized
leakage profile.

These three examples span the rate landscape. The Bloom-based
cipher map sits at the high-$\fprate$, zero-$\fnrate$ corner of the
perfect-filter line of chapter 6; the equality and ordering cipher
maps sit at the zero-zero corner but pay their confidentiality cost
through structural leakage rather than rate. The chapter's main
work is to show that the framework's $\fprate, \fnrate$ machinery
gives a quantitative grip on the first case, and that the framework's
PPV/NPV machinery (chapter 7 §7.4) gives a confidentiality measure
that includes the second.


% Section 17.2
% §17.2 The Bernoulli Framework as Error Model
%
% Monograph register: cipher-map rates ARE Bernoulli rates; per-query
% channel matrices; chained composition is matrix multiplication
% (thm:cipher-composition); perfect-filter geometry carries over;
% two-stage encrypted-classifier worked example. Tightened. Notes the
% trapdoor set as a Bernoulli<T> instance (salvaged ex-10 point, inline).
% Cross-refs sec:noisy-bit, sec:bernoulli-bool-formal, sec:kronecker,
% thm:kronecker, ax:element-indep, sec:perfect-filter-line,
% def:perfect-filter-line, sec:n-ary-channels, def:n-ary-channel,
% thm:n-ary-kronecker, def:cipher-map, sec:bernoulli-of-composite.
% DEFINED: sec:bernoulli-error-model, thm:cipher-composition.

\section{The Bernoulli Framework as Error Model}
\label{sec:bernoulli-error-model}

\Cref{def:cipher-map} stated the abstraction in terms of a
false-positive rate $\fprate$ and a false-negative rate $\fnrate$. That
notation was not cosmetic: the rates are the same objects chapters 1 and
4 introduced for the Bernoulli atom, and a chained cipher-map composition
obeys the same laws as a chained noisy-classifier composition. This
section makes the equivalence operational, so the algebraic identities
the framework proved earlier carry over to cipher maps unchanged.

\subsection*{Cipher-map rates are Bernoulli rates}

A Bernoulli predicate has per-query behavior summarized by $\fprate =
\Prob{\tilde A(x) = 1 \Given x \notin A}$ and $\fnrate = \Prob{\tilde
A(x) = 0 \Given x \in A}$, and \Cref{sec:bernoulli-bool-formal} showed
that the two rates with the input prevalence specify the joint
distribution of any finite query sequence. A cipher map $\tilde f$ for a
Boolean predicate $f$ has the same per-query behavior. The client's only
observable is the decrypted bit $\mathsf{Dec}(\tilde f(\mathsf{Enc}(x)))
\in \{0, 1\}$, and that bit is Bernoulli conditional on the plaintext
class:
\[
  \fprate = \Prob{\mathsf{Dec}(\tilde f(\mathsf{Enc}(x))) = 1 \Given f(x) = 0},
  \quad
  \fnrate = \Prob{\mathsf{Dec}(\tilde f(\mathsf{Enc}(x))) = 0 \Given f(x) = 1},
\]
the probability taken over $\tilde f$'s internal randomness (hash
choice, ciphertext noise, key schedule). Encryption and decryption drop
out because $\mathsf{Dec} \circ \mathsf{Enc}$ is the identity on
correctly decrypted outputs, leaving a Bernoulli atom with the cipher
map's error profile as its $(\fprate, \fnrate)$. Every theorem of
chapters 1 through 15 that took $(\fprate, \fnrate)$ as input applies
when the rates come from a cipher map; the framework is indifferent to
where the error originates and tracks only how it propagates. A trapdoor
set is in this sense a $\Bernoulli{T}$ over its plaintext type $T$
(\Cref{sec:bernoulli-of-composite}), the encryption layer changing the
representation but not the Bernoulli-square point.

\subsection*{Per-query channel matrices}

The two-rate summary is convenient but limited. A Boolean cipher map is
fully captured by a $2 \times 2$ channel matrix
\[
  Q \;=\;
  \begin{pmatrix}
    1 - \fprate & \fprate \\
    \fnrate & 1 - \fnrate
  \end{pmatrix},
\]
rows indexed by the true plaintext class, columns by the observed
output; the Bloom-based map has $\fnrate = 0$, a zero in the lower-left.
For non-Boolean cipher maps (encrypted classifiers, ranked retrieval)
the $2 \times 2$ matrix generalizes to the $K \times L$ channel matrix
of \Cref{def:n-ary-channel}, with $Q_{k,\ell} = \Prob{\mathsf{Dec}(\tilde
f(\mathsf{Enc}(x))) = \ell \Given f(x) = k}$. The channel-matrix
representation captures the full per-query behavior in one object that
composes by multiplication, and it makes contact with the
information-theoretic machinery of chapter 8 that
\Cref{sec:maximizing-confidentiality} needs to quantify confidentiality.

\subsection*{Chained composition is matrix multiplication}

Suppose the client computes $h = f \circ g$ on encrypted data, $X
\xrightarrow{g} Y \xrightarrow{f} Z$: a cipher map $\tilde g$ returns an
intermediate $\mathsf{Enc}(Y)$, then $\tilde f$ computes on it and
returns $\mathsf{Enc}(Z)$. The error compounds. If $\tilde g$ produces a
wrong $y' \ne g(x)$, then $\tilde f$ is asked to compute $f(y')$, and
even an errorless $\tilde f$ returns $f(y') \ne f(g(x))$: the stage-1
error propagates through stage 2. This is the situation the chapter 4
Kronecker theorem (\Cref{thm:kronecker}) and its n-ary generalization
(\Cref{thm:n-ary-kronecker}) were built to analyze, and the
factorization is matrix multiplication.

\begin{theorem}[Cipher-map composition]
\label{thm:cipher-composition}
Let $\tilde g$ and $\tilde f$ be cipher maps realizing $g : X \to Y$ and
$f : Y \to Z$ with per-query channel matrices $Q_g \in [0,1]^{|X| \times
|Y|}$ and $Q_f \in [0,1]^{|Y| \times |Z|}$. If \Cref{ax:element-indep}
(Axiom 1) holds, so the two stages' errors are mutually independent
given the inputs, then the composed cipher map $\tilde f \circ \tilde g$
realizes $h = f \circ g$ with channel matrix
\[
  Q_h \;=\; Q_g \, Q_f,
\]
ordinary matrix multiplication, $Q_h \in [0,1]^{|X| \times |Z|}$.
\end{theorem}

The proof reuses the chapter 4 argument: $(Q_g Q_f)_{x,z}$ sums over
intermediate $y$ the probability that stage 1 produces $y$ times the
probability that stage 2 then produces $z$, the conditional independence
being Axiom 1 applied to the two-stage sequence. For Boolean maps the
theorem is a $2 \times 2$ product; the same result governs any finite
chain by induction, the composed matrix being the product of the
per-stage matrices. A consequence: chained cipher-map error grows
predictably, the composed Boolean rates bounded by chapter 3's
$1 - \prod_i (1 - \eta_i)$ on each error type. This is Property 4 of the
cipher-map design in the trapdoor-computing literature, the framework's
composition law and the cryptographic side's requirement being the same
property.

\subsection*{The geometry carries over}

The perfect-filter line of \Cref{def:perfect-filter-line} classifies
approximate sets by their $(\fprate, \fnrate)$ position, and the same
classification applies to cipher maps. A cipher map with $\fnrate = 0$ is
\emph{safe}: it never misses a real positive, the client losing no real
results, at the cost of $\fprate > 0$. The Bloom-based map sits on this
line, as does any cipher map built by encrypting a perfect-filter-line
approximate set. A cipher map with $\fprate = 0$ sits on the opposite
axis, \emph{exact when positive} but missing some real positives;
it appears in randomized-encoding maps that occasionally reject valid
plaintexts. The classification is geometric and applies to cipher maps as
a special case of approximate functions.

\subsection*{Worked example: two-stage encrypted classifier}

A representative composition appears in encrypted retrieval. The client
retrieves and ranks documents matching an encrypted query in two stages.
\emph{Stage 1, encrypted membership}: the server holds an encrypted Bloom
filter on a candidate set $A$, returns an encrypted membership bit at
$\fprate_1 = 0.01$, $\fnrate_1 = 0$, so
\[
  Q_g = \begin{pmatrix} 0.99 & 0.01 \\ 0 & 1 \end{pmatrix}.
\]
\emph{Stage 2, exact ranking on positives}: an FHE-based ranking
introduces no error, $Q_f = I$. By \Cref{thm:cipher-composition}, $Q_h =
Q_g I = Q_g$: the composed map inherits $\fprate = 0.01$, $\fnrate = 0$.
Had stage 2 carried its own FPR-only error $\fprate_2 = 0.005$, the
product would give composed $\fprate = 0.99 \cdot 0.005 + 0.01 \cdot 1
\approx 0.015$, matching $1 - (1 - 0.01)(1 - 0.005) \approx 0.015$, with
$\fnrate$ staying zero because both stages are one-sided. The framework's
composition law applies to cipher maps with no modification. The next
section unpacks the operational consequences: when $\fprate, \fnrate$ are
also confidentiality parameters, the same matrix algebra governs how
confidentiality accumulates across composed queries.


% Section 17.3
% §17.3 The Trapdoor Computing Paradigm
%
% Monograph register: the paradigm (encrypt, compute on ciphertext, only
% key-holder decrypts), two-party server/client setup, confidentiality vs
% correctness on the perfect-filter line, Kronecker leakage over m
% queries, adaptive adversaries with a worked example.
% FIX: (1) notation collision in the worked example (the draft set A in
% {0,1}^n then wrote "n = |A|"); use N = |A| for dataset size throughout.
% (2) leakage-rate error (review M1): the draft claimed the adaptive
% adversary extracts O(log N) bits PER QUERY, which does not follow
% (seed-probing is O(1)/query; record-localization takes O(log N) queries
% TOTAL). Corrected: per-query extraction is O(1); the real gap is that a
% bounded budget of O(log N) crafted queries breaks the scheme outright,
% vs the iid attacker's O(fpr) bits/query that never recovers the seed.
% Cross-refs sec:kronecker, thm:kronecker, ax:element-indep,
% def:perfect-filter-line, sec:adversarial-inputs, def:cipher-map,
% sec:bernoulli-error-model, sec:fhe-sidebar, sec:maximizing-confidentiality.
% DEFINED: sec:trapdoor-paradigm.

\section{The Trapdoor Computing Paradigm}
\label{sec:trapdoor-paradigm}

The cipher map of \Cref{def:cipher-map} is the mathematical object; the
trapdoor computing paradigm is the operational setting in which it earns
its keep. The setting has three commitments: the data is encrypted before
it leaves the user, the computation runs entirely on the encrypted form,
and the result can be read only by a holder of a short secret key, the
trapdoor. The cipher map carries out the computation and the trapdoor
unlocks the answer; everything else, the plaintext records, the queries
that select among them, the intermediate results, stays hidden from any
party without the key.

\subsection*{The paradigm}

The paradigm predates the modern encryption-on-the-cloud framing.
Goldreich's work on software protection in the late 1980s asked how a
program could run on an untrusted machine without revealing what it
computed; the same question, in other language, drives oblivious RAM,
secure multiparty computation, and private information retrieval. What
unites them is the conviction that \emph{compute} and \emph{know} can be
separated: a machine can do useful work on data whose meaning it cannot
see. Trapdoor computing names the separation. The system has two
distinguished resources, an encrypted dataset and a trapdoor key, both
needed to read anything meaningful: the dataset without the key is opaque,
the key without the dataset useless. Computation is the operation of
cipher maps on the encrypted dataset, producing encrypted outputs the key
holder can decrypt and no one else can. The Bernoulli framework enters at
decryption: because the cipher map is approximate, the decrypted output
is a Bernoulli observation rather than a deterministic fact, and the
chapter 4 axioms govern how it behaves across queries. The cryptographic
side sets up the problem (compute without revealing); the Bernoulli side
tracks the consequence (read with controlled error).

\subsection*{Two-party setup}

The cleanest instance has a \emph{server} holding the encrypted dataset
and running the cipher maps, and a \emph{client} holding the trapdoor and
issuing queries. The server sees ciphertexts in and ciphertexts out,
nothing else. The encryption keeps the query private; the cipher map's
bounded error keeps the dataset partially hidden even from a client who
probes many times. The standard threat model is the
\emph{honest-but-curious} server, which follows the protocol but records
the encrypted traffic for later analysis; bounding what that analysis can
recover is the Bernoulli framework's job. A malicious server that deviates
from the protocol needs cryptographic integrity checks beyond this
chapter. The client's recovery quality is set by the $(\fprate, \fnrate)$
profile: an $\fnrate = 0$ map loses no real positives but admits some
false ones, an $\fprate = 0$ map trusts every positive but misses some
real ones.

\subsection*{Confidentiality versus correctness}

The trade-off that distinguishes the paradigm from cleartext computation
is between confidentiality and correctness. Tight error rates make the
client's view faithful but turn each near-deterministic answer into a
side channel: a server watching many such queries ties encrypted queries
to the records they match and reconstructs the plaintext distribution.
Loose rates make each query noisier for the client but push the
encrypted-output distribution toward uniform across the dataset, defeating
frequency analysis. The $(\fprate, \fnrate)$ point is therefore a
confidentiality knob as well as a correctness knob, and the
perfect-filter line of \Cref{def:perfect-filter-line} reads naturally in
this light: an $\fnrate = 0$ map concentrates its confidentiality cost in
the single parameter $\fprate$ (every real positive returned, recall
preserved), an $\fprate = 0$ map in $\fnrate$ (every positive answer
correct, but real positives missed). The first suits recall-sensitive
retrieval, the second precision-sensitive contexts.
\Cref{sec:maximizing-confidentiality} returns with a quantitative measure.

\subsection*{Multiple queries: leakage accumulates}

A single query reveals little; a million can reveal much. Under
\Cref{ax:element-indep} (Axiom 1) the error events of $m$ independent
queries are mutually independent, and the joint distribution factors as
the Kronecker product of $m$ copies of the per-query channel matrix
(\Cref{thm:kronecker}). The same algebra describes the trapdoor stream,
and any single-query confidentiality bound lifts through tensor algebra to
the stream. The independence cuts both ways: the client's recovery
sharpens as $m^{1/2}$ by standard concentration, and the server's
information about the plaintext grows linearly in $m$ for any map with
positive single-query leakage. The two scalings are the same
data-processing inequality read from opposite sides, the Kronecker
structure making the duality exact. A system with a finite query budget
can be analyzed by applying the per-query bounds to the $m$-fold Kronecker
product; an unbounded budget needs a different approach, either rotating
the trapdoor key to reset accumulated leakage or designing for the
short-horizon regime where asymptotic confidentiality has not yet decayed
to zero.

\subsection*{Adaptive adversaries}

The independence the Kronecker analysis depends on fails when the stream
is adaptive: the adversary chooses each query from the encrypted outputs
seen so far. This is exactly the adversarial-input setting of
\Cref{sec:adversarial-inputs}, and the trapdoor paradigm inherits its
failure modes, with the encrypted outputs (not plaintext ones) driving
the strategy. The bounds shift from average-case (iid Kronecker products)
to worst-case: a map at $\fprate = 0.01$ on iid queries may reach
$\fprate$ near $1$ on a query crafted from the preceding outputs. The
response, as in chapter 16, is worst-case accounting: bound the maximum
plaintext bits any adversary extracts per query and budget around that.

Consider a trapdoor membership map on a plaintext set $A$ of $N = |A|$
records, with $\fprate = 2^{-k}$ and $\fnrate = 0$, realized by a
Bloom-style construction with $k$ keyed hash functions whose seed is the
trapdoor. Under iid queries the behavior is benign: each random query has
independent probability $2^{-k}$ of a false positive, and after $m$
queries the server has gained $O(m)$ bits about $A$, spread uniformly
enough that no record is pinned down. Under an adaptive adversary the
picture changes, not in the per-query rate but in what a bounded budget
of crafted queries buys. The $\fprate = 2^{-k}$ comes from $k$ hash bits
matching at once; an adversary crafts queries that probe the bits
individually, each extracting about one bit of the seed, so after $O(k)$
such queries the seed is recovered and the entire dataset decrypts, and a
binary search over the records localizes a chosen record's membership in
$O(\log N)$ queries. The per-query extraction stays $O(1)$ bits
throughout; the danger is that this bounded budget, logarithmic in the
dataset size, breaks the scheme outright, where the iid attacker gains
only $O(\fprate)$ bits per query and never recovers the seed. The gap
between total compromise in a logarithmic number of crafted queries and
an $O(\fprate)$-bit-per-query trickle is the gap between worst-case and
average-case, and the fix, as chapter 16 prescribed, is a cipher map no
bounded budget of crafted queries can break. Cryptographic hash families
(pseudorandom functions) achieve this: a pseudorandom seed cannot be
recovered by probing, so the worst-case attacker gains only one bit of
best-effort guessing per query, matching the average case to a constant;
the FHE sidebar that follows gives an even stronger guarantee, worst-case
leakage exactly zero. Both
the worst-case story (chapter 16) and the average-case story (chapters 4
through 14) are needed in the trapdoor setting: the average case for
sizing under benign workloads, the worst case for security against
adversarial ones.


% Section 17.4
% §17.4 Sidebar: Fully Homomorphic Encryption
%
% Cross-reference resolution status:
%
% DEFINED here:
%   \label{sec:fhe-sidebar}                   -- this section
%
% ALREADY DEFINED (chapter 5):
%   \Cref{sec:bloom-from-scratch}             -- chapter 5 §5.2 Bloom filter
%
% ALREADY DEFINED (chapter 14):
%   \Cref{ch:optimal-bernoulli-sets}          -- chapter 14 BHF construction
%
% INTERNAL refs (chapter 17):
%   \Cref{def:cipher-map}                     -- §17.1
%   \Cref{sec:trapdoor-paradigm}              -- §17.3
%
% §17.4 is the FHE sidebar. Five items: (1) what FHE is, Gentry
% (2009); (2) cipher maps as a weaker, fprate/fnrate abstraction;
% (3) FHE cost (huge ciphertexts, slow operations); (4) when to use
% each (FHE for absolute correctness; cipher maps for tolerable
% error plus performance); (5) concrete comparison (Bloom cipher map
% at epsilon=0.01 vs FHE equivalent, ~10^4 space factor).

\section{Sidebar: Fully Homomorphic Encryption}
\label{sec:fhe-sidebar}

The previous sections built up cipher maps as approximate functions
on encrypted data. The approximate qualifier is a deliberate
weakening: a cipher map accepts bounded $\fprate$ and $\fnrate$ as
the price for being practical. A natural question is what the
unweakened version looks like. The answer is fully homomorphic
encryption, the cryptographic abstraction that promises zero-error
computation on ciphertexts. This sidebar steps off the main thread
of the chapter to lay out what FHE is, where it sits relative to
cipher maps, and when each abstraction is the right tool for the
job.

\subsection*{What fully homomorphic encryption is}

Fully homomorphic encryption (FHE) is an encryption scheme in which
the ciphertexts support exact arithmetic. Given $\mathsf{Enc}(x)$
and $\mathsf{Enc}(y)$, an FHE scheme provides operations
$\mathsf{Enc}(x) \oplus \mathsf{Enc}(y) = \mathsf{Enc}(x + y)$ and
$\mathsf{Enc}(x) \otimes \mathsf{Enc}(y) = \mathsf{Enc}(x \cdot y)$,
preserved exactly through any depth of composition. Because
arbitrary Boolean and arithmetic circuits decompose into additions
and multiplications, FHE supports the exact evaluation of any
computable function on encrypted inputs: a server can compute $\tilde
f(\mathsf{Enc}(x_1), \ldots, \mathsf{Enc}(x_n))$ and return
$\mathsf{Enc}(f(x_1, \ldots, x_n))$ with zero error in the
plaintext output.

Gentry's 2009 thesis gave the first viable construction of an FHE
scheme, resolving a question that had stood open since Rivest,
Adleman, and Dertouzos posed it in 1978. Earlier schemes had been
\emph{partially} homomorphic, supporting either addition (Paillier)
or multiplication (RSA, ElGamal) but not both, and so could not
evaluate arbitrary circuits. Gentry's construction allowed both
operations and showed how to control the ciphertext noise that
accumulates with depth. Subsequent schemes (BFV, BGV, CKKS, TFHE)
have refined the construction, traded off different cost dimensions,
and made FHE practical for a growing class of applications.

For the present chapter the relevant fact is that FHE realizes a
cipher map with $\fprate = \fnrate = 0$. The client recovers the
exact plaintext output of $f$, with no error from the cipher
construction itself. Whatever uncertainty the client has about the
plaintext after decrypting an FHE result is uncertainty that was
already present in the inputs, not uncertainty introduced by the
computation.

\subsection*{Cipher maps as a weaker abstraction}

A cipher map in the sense of \Cref{def:cipher-map} is the
strictly weaker abstraction. The cipher map admits a positive
$\fprate$ and $\fnrate$, and the framework's job is to bound the
error and propagate it through compositions. What the cipher map
buys in exchange for accepting error is room to operate. The
designer of a cipher map can choose representations (hash families,
sketch data structures, Bloom filters) that are orders of magnitude
smaller and faster than the polynomial-arithmetic representations
FHE requires. The Bernoulli framework's space-accuracy duality,
which has been the backbone of chapters 5 through 15, applies
directly: a cipher map for set membership at $\fprate = 2^{-20}$
uses about $20 + \mu$ bits per element by the optimal construction
of \Cref{ch:optimal-bernoulli-sets}, while the corresponding FHE
realization (an exact encrypted set) uses kilobytes per element.

The two abstractions therefore occupy different rungs on the same
ladder. FHE is the strongest abstraction the cryptographer can
deliver and the most expensive to use. Cipher maps are weaker (they
allow error) but cheaper in proportion to how much error the
application tolerates. The framework's role is to make the trade-off
quantitative: a cipher map with $\fprate = \epsilon$ uses about
$-\log_2 \epsilon$ bits per element of useful information capacity,
the cipher-map equivalent of FHE's exact-but-expensive guarantee.

\subsection*{The cost of FHE}

The price of FHE's zero-error guarantee is steep in three dimensions.

\textbf{Ciphertext size.} An FHE ciphertext encrypts a single
plaintext bit but occupies many kilobytes of memory. Early schemes
ran in the megabyte-per-bit range; modern BFV, BGV, and CKKS
ciphertexts pack tens to hundreds of plaintext bits into one
ciphertext, but still cost on the order of kilobytes per encrypted
bit at typical security parameters (128 to 256 bits of security
against quantum or classical attack). For a database that would
fit comfortably in a few gigabytes of plaintext, the FHE
representation runs to terabytes.

\textbf{Operation latency.} Each homomorphic addition or
multiplication takes milliseconds on a modern processor; deeper
circuits compound the cost. A simple lookup that takes nanoseconds
on plaintext data can take seconds or minutes under FHE, even on
the fastest available implementations. The slowdown is roughly six
to nine orders of magnitude relative to the plaintext operation,
depending on the scheme and the circuit depth.

\textbf{Bootstrapping.} The ciphertext noise that accumulates with
each homomorphic operation must periodically be reset; the
\emph{bootstrapping} operation that performs this reset is by far
the most expensive primitive in FHE, taking seconds even in the
fastest current implementations. Practical FHE workloads either
budget bootstrapping carefully (running only depth-bounded circuits
between resets) or accept the bootstrapping cost as the dominant
runtime component.

The combined cost is why FHE has not displaced plaintext computation
even for security-critical applications. The cryptographic community
has made enormous progress since 2009 (the cost has dropped by many
orders of magnitude), and FHE is now practical for narrow workloads
(simple aggregations, fixed-circuit medical queries, contract
evaluations on small data). But for the broad class of workloads
that cipher maps address (approximate set membership, ranked
retrieval, anonymity-preserving analytics over large datasets), FHE
remains too slow and too large.

\subsection*{When to use which}

The choice between FHE and a cipher map is the choice between exact
and approximate. Two questions decide it.

\textit{Does the application tolerate error?} A contract evaluation
on encrypted inputs is exact by definition: a single bit-flip in
the output changes the meaning of the contract. A medical record
lookup that returns whether a patient has a specific diagnosis is
exact in the sense that a false positive could trigger a wrongful
treatment. FHE is the right choice for these correctness-critical
queries; cipher maps are not. A search query that returns the top
ten matching documents from a million-document corpus tolerates a
small false-positive rate (a few irrelevant documents in the top
ten are an inconvenience, not a catastrophe). A trapdoor cipher
map at $\fprate = 0.01$ is the right choice here.

\textit{Does the application need to scale?} A handful of exact
queries per day on a small encrypted dataset can run under FHE
without difficulty: even at minutes per query, the throughput is
adequate. A search index serving thousands of queries per second
on a multi-gigabyte dataset cannot run under FHE at any realistic
hardware budget. Cipher maps scale because they live in the same
asymptotic regime as plaintext data structures: a Bloom-filter
cipher map serves queries in constant time and uses space
proportional to the dataset size, the same complexity as a
cleartext Bloom filter.

In practice the two abstractions often combine. A hybrid system
might use a cipher map as the first-stage filter (rejecting most
non-matches quickly with bounded error) and FHE as the second-stage
verifier (exactly checking the cipher map's positive results). The
cipher map keeps the per-query cost low; the FHE stage cleans up
the false positives. The framework's composition theorem of
\Cref{thm:cipher-composition} applies to this hybrid: the
overall $(\fprate, \fnrate)$ is the matrix product of the two
stages, and the FHE stage's exact channel (identity matrix) drops
the false-positive rate to whatever the cipher map can verify.

\subsection*{A concrete comparison}

A worked example fixes the cost ratio. Take a set membership query
on a plaintext set $A$ of $|A| = 10^6$ elements, each represented
by a 64-bit integer.

\textbf{Bloom-filter cipher map at $\fprate = 0.01$.} The optimal
Bloom-filter construction of chapter 5 uses about $-\log_2(0.01)
\approx 6.64$ bits per element on the underlying filter, plus a
small constant overhead for hash-family parameters. The trapdoor
adds a keyed seed (a few dozen bytes total) but does not change the
per-element cost asymptotically. The total cipher-map representation
is about $10^6 \times 6.64 \approx 6.64 \times 10^6$ bits, or
roughly 830 kilobytes. Each membership query runs in constant time
(a few hash computations and bit lookups), with no bootstrapping or
ciphertext expansion overhead.

\textbf{FHE-encrypted exact set.} An FHE realization of the same
set membership query encrypts each element under a scheme like BFV
or BGV. At a security parameter of 128 bits, each ciphertext
occupies about $10^4$ to $10^5$ bits, or roughly 1 to 10 kilobytes
per encrypted element. The total representation is about $10^6
\times 10^4 = 10^{10}$ bits, or roughly 1 to 10 gigabytes for the
same 1 million elements. Each membership query under FHE requires
homomorphic comparisons against each ciphertext or against a
homomorphic hash, taking seconds in either case.

\textbf{The ratio.} The cipher map costs about 830 kilobytes; the
FHE version costs about 1 to 10 gigabytes. The space ratio is on
the order of $10^4$, with the cipher map four orders of magnitude
smaller. The runtime ratio is roughly the same: nanoseconds for the
cipher map, seconds for FHE, giving about $10^9$ in favor of the
cipher map. The four-orders-of-magnitude space gap and nine-orders
runtime gap are the price the cipher map pays in correctness (a
$\fprate$ of 1 percent) for performance practical enough to serve a
production search engine.

This is the central trade-off the trapdoor computing paradigm
exploits. FHE remains the abstraction of choice when error is
intolerable. Cipher maps fill the much larger space of workloads
where bounded error is acceptable, and the framework's job is to
keep the error rates honest as the cipher maps compose. The next
section returns to the cipher-map setting and asks how Boolean
combinations of cipher maps behave under the framework's
composition rules.


% Section 17.5
% §17.5 Boolean Algebra Over Trapdoor Sets
%
% Cross-reference resolution status:
%
% DEFINED here:
%   \label{sec:boolean-trapdoor}              -- this section
%   \label{def:trapdoor-set}                  -- trapdoor set definition
%
% ALREADY DEFINED (chapter 3):
%   \Cref{thm:composition}                    -- chapter 3 §3.3 composition theorem
%
% ALREADY DEFINED (chapter 5):
%   \Cref{sec:bloom-from-scratch}             -- chapter 5 §5.2 Bloom filter
%
% ALREADY DEFINED (chapter 6):
%   \Cref{sec:cancellation}                   -- chapter 6 §6.3 cancellation analysis
%   \Cref{def:perfect-filter-line}            -- chapter 6 §6.4 definition
%   \Cref{prop:complement-axis-swap}          -- chapter 6 §6.3 complement axis swap
%
% INTERNAL refs (chapter 17):
%   \Cref{def:cipher-map}                     -- §17.1
%   \Cref{sec:trapdoor-paradigm}              -- §17.3
%   \Cref{thm:cipher-composition}             -- §17.2
%
% §17.5 introduces Boolean algebra over trapdoor sets. Five items:
% (1) trapdoor sets: encrypted set representations needing key for
% membership; Bloom is example; (2) Boolean ops (union, intersection,
% complement) on trapdoor representations; (3) composition theorem
% (chapter 3) applies to trapdoor-set composition; (4) cancellation
% (chapter 6 §6.3) lets trapdoor composition escape perfect-filter
% line; (5) worked example: query x in A AND x in B on two trapdoor
% sets, derive overall fprate.

\section{Boolean Algebra Over Trapdoor Sets}
\label{sec:boolean-trapdoor}

The trapdoor cipher maps of \Cref{sec:trapdoor-paradigm} are
useful in isolation: a single encrypted Bloom filter answers
membership queries with controlled error. They become more useful
when several are combined. A common request in encrypted-search
applications is the conjunctive query: \emph{is $x$ in $A$ AND in
$B$?}, asked of two trapdoor sets held on the same server. The
answer requires Boolean reasoning over encrypted representations,
not over plaintext sets. This section asks how Boolean algebra
behaves when its operands are trapdoor sets, and finds that the
framework's composition rules continue to apply, with the same
geometric subtleties that chapter 6 raised for plaintext approximate
sets.

\subsection*{Trapdoor sets}

A trapdoor set is an encrypted representation of a plaintext set
whose membership query requires a secret key.

\begin{definition}[Trapdoor set]
\label{def:trapdoor-set}
Let $A \subseteq X$ be a plaintext set and let $k$ be a secret key.
A \emph{trapdoor set} for $A$ under $k$, denoted $\tilde A$, is an
encrypted data structure together with an associated membership
cipher map
\[
  \tilde m_A : \mathsf{Enc}_k(X) \to \mathsf{Enc}_k(\{0, 1\})
\]
such that, for $x \in X$, $\mathsf{Dec}_k(\tilde m_A(\mathsf{Enc}_k(x))) =
\mathbf{1}[x \in A]$ with bounded error rates $\fprate_A, \fnrate_A$ in
the sense of \Cref{def:cipher-map}. Without the key $k$, the
membership query is computationally infeasible; with the key, it is
as cheap as a plaintext set lookup.
\end{definition}

The Bloom filter of chapter 5 \Cref{sec:bloom-from-scratch} is the
canonical example. A server holds an encrypted bit array
representing the Bloom filter's state, plus encrypted hash family
parameters that depend on the key. The client encrypts a query
element under the key, sends the ciphertext to the server, and the
server returns the encrypted membership bit. The cipher map has
$\fprate_A = (1 - e^{-kn/m})^k$ (the standard Bloom-filter formula
with $k$ hash functions, $n$ inserted elements, and $m$ bits), and
$\fnrate_A = 0$ from the one-sided error of the construction. The
trapdoor set sits at the point $(\fprate_A, 0)$ on the perfect-filter
line of \Cref{def:perfect-filter-line}.

Other trapdoor-set constructions sit at different points on the
rate plane. A perfect-hash-filter trapdoor set (chapter 15) sits
at $(2^{-r}, 0)$ for an $r$-bit fingerprint. A counting trapdoor
set built on the same hash family can support deletions at the
cost of a small $\fnrate_A$. The space of constructions is large;
what they share is the structural property of \Cref{def:trapdoor-set}:
encrypted representation, key-dependent membership query, bounded
Bernoulli error.

\subsection*{Boolean operations on trapdoor sets}

The three primitive Boolean operations on sets are union $A \cup B$,
intersection $A \cap B$, and complement $\complement A$.
Each lifts to the trapdoor setting, but the lift requires the server
to compute on encrypted representations rather than plaintext.

\textbf{Union.} A trapdoor set $\tilde{A \cup B}$ realizes the
plaintext union $A \cup B$ in the sense that
\[
  \mathsf{Dec}_k(\tilde m_{A \cup B}(\mathsf{Enc}_k(x))) =
  \mathbf{1}[x \in A \cup B] \;=\; \mathbf{1}[x \in A] \lor \mathbf{1}[x \in B],
\]
up to Bernoulli error. For Bloom-filter trapdoor sets the union can
be computed by a bitwise OR of the encrypted bit arrays, provided
the encryption scheme is additively homomorphic over the bits. The
server returns the encrypted bit obtained by querying both arrays
and OR-ing the results, and the client decrypts to read the union
membership.

\textbf{Intersection.} The intersection $\tilde{A \cap B}$ is the
analogous lift of $A \cap B$. The server queries both trapdoor sets
and AND-s the encrypted membership bits, returning an encrypted
result that decrypts to $\mathbf{1}[x \in A] \land \mathbf{1}[x \in B]$.
Unlike union, intersection on Bloom-filter trapdoor sets cannot
generally be computed by a bitwise operation on the underlying
arrays (the bitwise AND of two Bloom-filter arrays is not the Bloom
filter of the intersection, a well-known pitfall in the cleartext
case). Instead the server runs both membership queries separately
and combines the encrypted output bits with a homomorphic AND.

\textbf{Complement.} The complement $\tilde{\complement A}$ is
computed by complementing the encrypted membership bit:
$\mathsf{Dec}_k(\tilde m_{\complement A}(\mathsf{Enc}_k(x))) =
1 - \mathbf{1}[x \in A]$ up to Bernoulli error. The server flips
the returned bit (or runs the cipher map in complement mode). The
operation is cheap; what changes, as \Cref{prop:complement-axis-swap}
warned, is the position of the result on the perfect-filter line.

These three primitives suffice to build any finite Boolean
combination by composition. The interesting questions concern how
the operations' error rates compose, and how confidentiality
propagates when the server has access to the operations' encrypted
outputs.

\subsection*{The composition theorem applies}

The error rates of a Boolean combination of trapdoor sets are the
error rates of the corresponding plaintext Boolean operation,
computed by the framework's composition theorem.
\Cref{thm:composition} stated the rule for plaintext approximate
sets: when an operation combines $k$ approximate operands whose
error events are mutually independent under Axiom 1, the resulting
overall error rate satisfies
\[
  \eta_{\text{total}} \;=\; 1 - \prod_{i=1}^{k} (1 - \eta_i),
\]
applied separately to the FPR and FNR channels.

The same formula applies to trapdoor-set composition without
modification. Two trapdoor sets $\tilde A, \tilde B$ with per-set
rates $(\fprate_A, 0)$ and $(\fprate_B, 0)$, combined under union,
yield a trapdoor set $\tilde{A \cup B}$ with rates
\[
  \fprate_{A \cup B} \;=\; 1 - (1 - \fprate_A)(1 - \fprate_B),
  \qquad
  \fnrate_{A \cup B} \;=\; 0.
\]
Combined under intersection the rates are
\[
  \fprate_{A \cap B} \;=\; \fprate_A \cdot \fprate_B,
  \qquad
  \fnrate_{A \cap B} \;=\; 0,
\]
because a query $x \notin A \cap B$ is a false positive only if it
is a false positive under both component membership tests, which
happens with probability $\fprate_A \cdot \fprate_B$ under
independence.

The conditional-independence hypothesis is the chapter 4 Axiom 1
applied to the trapdoor setting: the two trapdoor sets' cipher
maps use independent randomness (independent hash family seeds,
independent encryption noise) and so produce independent false-positive
events for any given query. This is automatic if the two trapdoor
keys are generated independently, and is the standard assumption in
encrypted-search constructions.

The same algebra propagates through deeper compositions. A
three-way conjunctive query $x \in A \cap B \cap C$ on trapdoor
sets at rates $\fprate_A, \fprate_B, \fprate_C$ has overall
$\fprate = \fprate_A \fprate_B \fprate_C$. A $k$-way conjunction
has $\fprate = \prod_i \fprate_i$, which decays exponentially in
$k$: deep conjunctive queries are highly precise even when the
individual trapdoor sets are loose. The framework's composition
law captures this strengthening exactly, and the trapdoor-set
geometry inherits it.

\subsection*{Cancellation escapes the perfect-filter line}

The composition formula above kept the operands on the perfect-filter
line: both unions and intersections of FNR-zero trapdoor sets stayed
FNR-zero. What chapter 6 \Cref{sec:cancellation} flagged is that
this clean behavior is not a general property of Boolean composition;
it depends on the operands sitting on the same axis. As soon as the
complement enters the picture (or as soon as one operand sits on the
opposite axis from another), the result can leave the perfect-filter
line entirely.

The cancellation phenomenon transfers directly to trapdoor sets.
Consider a trapdoor union $\tilde A \cup \tilde B$ at rates
$(\fprate_A, 0)$ and $(\fprate_B, 0)$, both on the FNR-zero axis.
Now take the complement of this union, $\tilde{\complement(A \cup B)}$.
By \Cref{prop:complement-axis-swap} the complement of an FNR-zero
trapdoor set is FPR-zero, so $\tilde{\complement(A \cup B)}$ sits
at the point $(0, \fprate_{A \cup B})$ on the FPR-zero axis. The
trapdoor system has moved off the original axis.

Compose further. Suppose the server is asked $\tilde{\complement(A
\cup B)} \cap \tilde C$, where $\tilde C$ is a third trapdoor set
back on the FNR-zero axis at rates $(\fprate_C, 0)$. The
intersection now mixes operands from different axes. By the rate
algebra,
\[
  \fprate_{\complement(A \cup B) \cap C} \;=\;
  \fprate_{\complement(A \cup B)} \cdot \fprate_C \;=\; 0,
\]
because one operand is FPR-zero. But the FNR is no longer zero,
because the FPR-zero operand carries a positive $\fnrate$:
\[
  \fnrate_{\complement(A \cup B) \cap C} \;=\;
  1 - (1 - \fprate_{A \cup B})(1 - \fnrate_C)
  \;=\; \fprate_{A \cup B},
\]
since $\fnrate_C = 0$. The composed trapdoor set has rates
$(0, \fprate_{A \cup B})$ and now sits on the FPR-zero axis: the
two complements and the cross-axis intersection have moved the
result through the rate plane in a manner the perfect-filter-line
classification cannot describe by a single one-sidedness flag.

More elaborate Boolean combinations can produce trapdoor sets with
both rates positive, fully off the perfect-filter line. This is
the same cancellation phenomenon chapter 6 documented for
plaintext approximate sets, lifted intact to the trapdoor setting.
The practical implication is that designing trapdoor-set Boolean
combinations requires the same care as designing plaintext Boolean
combinations: track the operand axes through the composition, and
expect the perfect-filter-line classification to dissolve as soon
as the composition mixes axes.

\subsection*{Worked example: conjunctive query}

A concrete example fixes the rate algebra. Take two trapdoor Bloom
filters $\tilde A$ and $\tilde B$ on plaintext sets $A, B \subseteq
X$, each configured at $\fprate = 0.01$ with $\fnrate = 0$. The
client wants to answer the query
\[
  q(x) \;=\; \mathbf{1}[x \in A] \;\land\; \mathbf{1}[x \in B] \;=\;
  \mathbf{1}[x \in A \cap B].
\]
The client encrypts $x$ under the shared key, sends $\mathsf{Enc}_k(x)$
to the server, and asks the server to compute the conjunctive
membership cipher map.

\textbf{Server computation.} The server queries both trapdoor sets,
producing encrypted bits $\tilde b_A = \tilde m_A(\mathsf{Enc}_k(x))$
and $\tilde b_B = \tilde m_B(\mathsf{Enc}_k(x))$. It combines them
under a homomorphic AND to produce $\tilde b = \tilde b_A \land
\tilde b_B$, and returns $\tilde b$ to the client. The client
decrypts: $\mathsf{Dec}_k(\tilde b)$ is a bit in $\{0, 1\}$ which
approximates $\mathbf{1}[x \in A \cap B]$ up to the composed
Bernoulli error.

\textbf{Rate analysis.} Apply the composition formula above. The
intersection of two FNR-zero trapdoor sets has $\fnrate = 0$ (no
real positive is lost, because both component queries return $1$
on true positives). The composed FPR is the product:
\[
  \fprate_{A \cap B} \;=\; \fprate_A \cdot \fprate_B \;=\; 0.01 \times 0.01
  \;=\; 10^{-4}.
\]
The conjunctive query has $\fprate = 10^{-4}$, two orders of
magnitude tighter than either component trapdoor set. The client
receives an answer that is correct on all true positives (because
$\fnrate = 0$) and rejects all but one in ten thousand false
positives.

\textbf{Interpretation.} The intersection sharpens precision
multiplicatively. Each additional trapdoor set in a $k$-way
conjunction divides the false-positive rate by another factor on
the order of $\fprate$, and a $k$-way conjunction has $\fprate
= \fprate^k$ (assuming equal-rate operands and independent
trapdoors). The framework's composition theorem captures this
sharpening exactly; the trapdoor setting inherits it without
qualification.

This is the most useful pattern in encrypted search built from
trapdoor sets. A query that asks for documents matching all of
several encrypted keywords becomes a conjunction of trapdoor-set
memberships, and the conjunction's precision sharpens exponentially
in the number of keywords. The same composition law that governs
plaintext Bloom-filter conjunctions governs the trapdoor lift:
matrix multiplication in the channel-matrix formulation of
\Cref{thm:cipher-composition}, or equivalently the product
formula above when the operands are FNR-zero. The next section turns
to how this rate-algebra story is also a confidentiality-algebra
story, and asks how much information about the plaintext sets the
server can recover from the conjunctive query stream.


% Section 17.6
% §17.6 Maximizing Confidentiality
%
% Monograph register: confidentiality as the 1-PPV posterior; the ch7 PPV
% formula IS the confidentiality formula; trading rate-budget for
% confidentiality on the perfect-filter line; worked example. Tightened.
% Folds salvaged ex-11 as rem:confidentiality-bits (information-theoretic
% leakage L = H(pi) - E[H(posterior)] in bits/query, complementing
% C = 1 - PPV). Cross-refs def:perfect-filter-line, def:ppv, def:npv,
% def:cipher-map, sec:trapdoor-paradigm, sec:fhe-sidebar. DEFINED:
% sec:maximizing-confidentiality, rem:confidentiality-bits.

\section{Maximizing Confidentiality}
\label{sec:maximizing-confidentiality}

\Cref{sec:trapdoor-paradigm} named a duality the framework has not yet
quantified: the cipher-map rates $\fprate$ and $\fnrate$ govern both
correctness and confidentiality. Looser rates make the server's output
less faithful to $f$ but also less informative to an adversary inspecting
it. The earlier sections treated correctness; this one treats
confidentiality, and the question is concrete: given a cipher map at fixed
$(\fprate, \fnrate)$ on a workload of known prevalence $\pi$, how much
does the output reveal about the plaintext? The chapter 7 PPV machinery
answers it with no further apparatus.

\subsection*{Confidentiality is a posterior}

An adversary watching the server's outputs faces a posterior-inference
problem: each output is a noisy report about the plaintext, and Bayes'
rule combines it with prior knowledge. For a single Boolean query the
target is
\[
  \Pr\bigl[\,x \in A \;\big\vert\; \tilde f(x) = \top\,\bigr],
\]
the probability that a reported $\top$ identifies a genuine member rather
than a false positive. Close to $1$, the report has pinned membership
down; close to the prior $\pi$, it has revealed nothing. The
confidentiality is the gap, which we take to be
\[
  C \;=\; 1 - \PPV,
\]
the share of $\top$-reports that are false positives. Confidentiality is
high when most $\top$-reports do not identify members and zero in the FHE
limit, where every report is correct and membership is fully exposed.

The posterior is exactly the positive predictive value of
\Cref{def:ppv}. Reading the chapter 7 formula in the cipher-map setting,
\[
  \PPV = \frac{\pi(1 - \fnrate)}{\pi(1 - \fnrate) + (1 - \pi)\fprate},
  \qquad
  C = 1 - \PPV = \frac{(1 - \pi)\fprate}{\pi(1 - \fnrate) + (1 - \pi)\fprate},
\]
the numerator $\pi(1-\fnrate)$ the rate of correct $\top$ on a member,
$(1-\pi)\fprate$ the rate of spurious $\top$ on a non-member, $C$ the
share that protects the plaintext. Zero false positives gives $C = 0$; an
unconditional $\top$ gives $C = 1 - \pi$. The negative case mirrors it,
the confidentiality of a $\bot$-report being $1 - \NPV$ (\Cref{def:npv}),
though in typical encrypted search $\pi$ is small and the $\top$-side
ambiguity dominates.

The $C = 1 - \PPV$ measure is the per-report complement; the
information-theoretic complement, in bits, is also available.

\begin{remark}[Confidentiality in bits per query]
\label{rem:confidentiality-bits}
Model the adversary's per-query posterior as Bernoulli at $\PPV$ for a
$\top$-report and at $1 - \NPV$ for a $\bot$-report. The leakage is the
mutual information between the plaintext bit and the report,
\[
  L \;=\; H(\pi) - \mathbb{E}_{\text{report}}\!\left[H(\text{posterior})\right]
  \;=\; H(\pi) - \bigl[\Pr(\top)\,H(\PPV) + \Pr(\bot)\,H(1 - \NPV)\bigr],
\]
with $H$ the binary entropy and $\Pr(\top) = \pi(1-\fnrate) +
(1-\pi)\fprate$. In the FHE limit $\fprate = \fnrate = 0$ the posteriors
are deterministic, $L = H(\pi)$, the full prior entropy leaks. A cipher
map with positive rates leaks a fraction of $H(\pi)$ per query, the
information-theoretic counterpart of the $C = 1 - \PPV$ posterior, and the
chapter 8 entropy machinery measures that fraction directly.
\end{remark}

\subsection*{Trading rate-budget for confidentiality}

On the perfect-filter line of \Cref{def:perfect-filter-line} the designer
faces a one-dimensional trade-off. An $\fnrate = 0$ map sits on the
$\fprate$-axis, and raising $\fprate$ raises $(1-\pi)\fprate$ and so
raises $C$ monotonically. Walking the line away from the origin trades
correctness for confidentiality. The trade is not free in the convenient
direction: chapter 5's space formula makes storage a decreasing function
of $\fprate$, so \emph{lower-storage cipher maps are more confidential},
one of the framework's small surprises, following entirely from the
chapter 7 posterior with no cryptographic argument. Storage couples to
accuracy through the chapter 8 entropy bound and accuracy to
confidentiality through the chapter 7 posterior, so one parameter sets
three coupled costs. Maximizing $C$ in the abstract is degenerate (walk
to $\fprate = 1$, $C = 1 - \pi$, the map useless), so the real problem
constrains the operational error: maximize $C$ subject to $\fprate \le
\fprate^*$ for an application-imposed cap, the optimum being the cap
itself.

\subsection*{A worked example}

Take prevalence $\pi = 0.01$ and rates $\fprate = 0.05$, $\fnrate = 0$.
The PPV formula gives
\[
  \PPV = \frac{0.01}{0.01 + 0.99 \cdot 0.05} = \frac{0.01}{0.0595}
  \approx 0.168,
  \qquad C \approx 0.832.
\]
When the server reports $x \in A$, the adversary's posterior assigns about
$16.8\%$ to genuine membership and $83.2\%$ to a false positive; one query
leaves the question far from settled. The FHE baseline reports with zero
error, $\PPV = 1$, $C = 0$: every report is correct and fully revealing.
The contrast is stark, the cipher map paying a $5\%$ false-positive rate
for an $83\%$ confidentiality margin, FHE paying a $10^4\times$ space and
$10^9\times$ time overhead for a zero margin. The example also shows the
workload dependence: holding $(\fprate, \fnrate) = (0.05, 0)$ but moving
to a balanced workload $\pi = 0.5$ gives $\PPV = 0.5/0.525 \approx 0.952$
and $C \approx 0.048$. The same map is confidential on a rare-positive
workload and revealing on a balanced one. The framework keeps the
parameters separate: the map carries $(\fprate, \fnrate)$, the workload
carries $\pi$, and $C(\fprate, \fnrate, \pi)$ is one function the designer
plots and reads. Chapter 7 did the algebra; chapter 17 inherits it.


% Section 17.7
% §17.7 Bridge
%
% Cross-reference resolution status:
%
% DEFINED here:
%   \label{sec:ch17-bridge}                   -- this section
%
% INTERNAL refs (chapter 17):
%   \Cref{sec:what-is-cipher-map}             -- §17.1
%   \Cref{sec:bernoulli-error-model}          -- §17.2
%   \Cref{sec:trapdoor-paradigm}              -- §17.3
%   \Cref{sec:fhe-sidebar}                    -- §17.4
%   \Cref{sec:boolean-trapdoor}               -- §17.5
%   \Cref{sec:maximizing-confidentiality}     -- §17.6
%   \Cref{sec:bib-notes-ch17}                 -- §17.8 (Task 9)
%
% RESOLVED (chapter labels for forward refs):
%   \Cref{ch:cipher-maps-trapdoor}            -- this chapter
%   Chapter 18 (encrypted search)             -- Plan 0 stub
%
% §17.7 closes chapter 17 with three items: (1) recap of the dual
% role of fprate/fnrate (correctness AND confidentiality);
% (2) chapter 18 forward-link to encrypted search as a specific
% cipher-map family; (3) signpost to §17.8 antecedents.

\section{Bridge}
\label{sec:ch17-bridge}

The chapter has built a single bridge in six pieces. The cipher-map
abstraction of \Cref{sec:what-is-cipher-map} reframed the
encryption setting as a pair of error rates rather than a pair of
keys: the server's job is to approximate $f$ on encrypted inputs,
and the bounded $\fprate$ and $\fnrate$ control how well it
succeeds. The error-model section
(\Cref{sec:bernoulli-error-model}) identified those rates with the
Bernoulli atom of chapter~4 and lifted chapter~4's Kronecker
composition law to chained cipher-map evaluation. The
trapdoor-paradigm section (\Cref{sec:trapdoor-paradigm}) introduced
the client-server split that makes the abstraction useful in
practice, and the FHE sidebar (\Cref{sec:fhe-sidebar}) marked the
zero-error baseline against which the cipher-map approach is
measured. Sections~\ref{sec:boolean-trapdoor}
and~\ref{sec:maximizing-confidentiality} pushed the framework
forward in two directions: Boolean composition over trapdoor sets
(chapter~3's composition theorem in cryptographic clothing) and
confidentiality as a chapter~7 PPV.

\subsection*{The dual role of the error rates}

The recurring thread is the dual role of $\fprate$ and $\fnrate$.
In Parts~I-IV the two rates were operational parameters: they
described how often the framework's structures returned wrong
answers, and the design problem was to keep them small subject to
storage and runtime budgets. In chapter~17 the same parameters
acquired a second job. The cipher-map false-positive rate is also
the share of the adversary's $\top$-observations that fail to
identify members of $A$, and the cipher-map false-negative rate is
also the share of $\bot$-observations that protect real members.
The two rates together fix the cipher map's PPV at workload
prevalence $\pi$, and the complementary $1 - \PPV$ is the
confidentiality margin.

The duality is not a coincidence. The Bernoulli framework
identifies a structure by its $(\fprate, \fnrate)$ point in the
Bernoulli square; chapter~6 read that point geometrically as a
position on the perfect-filter line; chapter~7 read the same point
predictively as PPV and NPV. Chapter~17 has read the point a third
time, cryptographically, as the trade-off between operational
fidelity and informational opacity. The three readings are the
same algebra. Designing a cipher map means choosing where to live
in the Bernoulli square, and that single choice now answers three
different questions at once.

\subsection*{Forward to chapter~18: encrypted search}

Chapter~18 takes the cipher-map machinery and specializes it to
one application: encrypted search. The setting fixes $f$ to a
membership-test family ``does record $r$ match query $q$?'' and
asks how a server that holds an encrypted index can answer such
queries faithfully without revealing $r$ or $q$ to whoever is
running the server. Bloom-filter cipher maps from
\Cref{sec:boolean-trapdoor} reappear as the simplest index
structure, and the chapter~14 BHF construction reappears as a
space-optimal alternative. The two-party setup of
\Cref{sec:trapdoor-paradigm} is still in force; the
$(\fprate, \fnrate)$ point still does triple duty as
correctness, geometry, and confidentiality.

What chapter~18 adds is the structure that the encrypted-search
literature has built around the membership-test case: oblivious
access patterns, side-channel considerations, multi-query leakage
bounds, and the relation between confidentiality $C$ and the
information-theoretic measures the encrypted-search literature
prefers. Chapter~17's algebraic machinery sets up the formal
language; chapter~18 fills in the application-specific structure.
The conceptual move from chapter~17 to chapter~18 is the same
move chapters~14 and~15 made together: chapter~17 is the abstract
algebra, chapter~18 is the concrete instance.

\subsection*{Antecedents}

The cipher-map abstraction and the trapdoor-computing paradigm
have antecedents in the cryptographic literature that the chapter
has named only in passing. The FHE-sidebar section pointed to
Gentry's 2009 construction as the zero-error baseline; the
trapdoor-paradigm section borrowed terminology from the
oblivious-computing tradition that chapter~16's bibliographic
notes already touched. The next section
(\Cref{sec:bib-notes-ch17}) attaches the chapter's prose to these
antecedents formally and gives the reader entry points to the
four literatures the chapter has been drawing on.


% Section 17.8
% §17.8 Bibliographic Notes
%
% Cross-reference resolution status:
%
% DEFINED here:
%   \label{sec:bib-notes-ch17}                -- this section
%
% INTERNAL refs (chapter 17):
%   \Cref{sec:what-is-cipher-map}             -- §17.1
%   \Cref{sec:fhe-sidebar}                    -- §17.4
%   \Cref{sec:boolean-trapdoor}               -- §17.5
%   \Cref{sec:maximizing-confidentiality}     -- §17.6
%
% Citations introduced here:
%   \cite{towell-cipher-maps}                 -- new bib entry
%   \cite{towell-boolean-trapdoor}            -- new bib entry
%   \cite{towell-maximizing-confidentiality}  -- new bib entry
%   \cite{gentry-2009}                        -- new bib entry
%   \cite{naor-yung-1989}                     -- new bib entry
%   \cite{goldreich-1996}                     -- already in bib from ch16
%
% Four paragraphs: cipher-maps manuscript; trapdoor-computing
% series (Boolean algebra over trapdoor sets, maximizing
% confidentiality, plus the algebraic cipher types thread);
% FHE literature (Gentry 2009); oblivious-computing tradition
% (Goldreich 1996, Naor-Yung 1989).

\section*{Bibliographic Notes}
\addcontentsline{toc}{section}{Bibliographic Notes}
\label{sec:bib-notes-ch17}

\paragraph{The cipher-maps manuscript.}
The cipher-map abstraction of \Cref{sec:what-is-cipher-map} is
developed at length in a companion manuscript,
\textcite{towell-cipher-maps}, ``Cipher maps: a unified framework
for oblivious function evaluation,'' which the present chapter
summarizes. The manuscript treats the abstraction in its own right
rather than as a special case of the Bernoulli framework, gives a
classification of cipher-map families by their channel-matrix
structure, and develops the cryptographic-side analysis (key
schedules, side-channel considerations, computational hardness
assumptions) that the present chapter has not entered. The
chapter's contribution beyond the manuscript is the identification
of cipher-map error rates with Bernoulli atoms and the lifting of
chapter~4's Kronecker composition to cipher-map composition; the
manuscript treats composition algebraically without anchoring it to
the Bernoulli square.

\paragraph{The trapdoor-computing series.}
The trapdoor-computing paradigm is the subject of an ongoing series
of which three drafts are directly relevant. The Boolean-algebra
section (\Cref{sec:boolean-trapdoor}) draws on
\textcite{towell-boolean-trapdoor}, ``Boolean algebra over trapdoor
sets,'' which gives the full algebraic treatment of union,
intersection, complement, and conditional operations over trapdoor
representations, including the cancellation phenomena the
present chapter alluded to. The maximizing-confidentiality section
(\Cref{sec:maximizing-confidentiality}) draws on
\textcite{towell-maximizing-confidentiality}, ``Maximizing
confidentiality in encrypted search,'' which develops the
PPV-as-confidentiality connection in the encrypted-search setting
and extends the analysis to information-theoretic measures
(observed-versus-maximum entropy ratios, compression-based entropy
estimation) the present chapter has not discussed. A third strand,
on algebraic cipher types, generalizes the chapter's set-and-map
treatment to algebraic data types and remains in active drafting.

\paragraph{Fully homomorphic encryption.}
The FHE sidebar of \Cref{sec:fhe-sidebar} stands on
\textcite{gentry-2009}, ``Fully Homomorphic Encryption Using Ideal
Lattices,'' which gave the first plausibly-secure fully
homomorphic scheme by combining a somewhat-homomorphic scheme with
a bootstrapping construction that refreshes ciphertexts faster
than the noise from homomorphic operations accumulates. The
scheme is conceptually clean and operationally heavy; modern
schemes (BFV, BGV, CKKS) reduce the per-operation cost by orders
of magnitude but inherit Gentry's bootstrapping idea and remain
several orders of magnitude slower than the cipher-map approach
that the present chapter develops. The cryptographic standard
references on FHE are surveys rather than primary sources, and
the reader interested in modern schemes should follow the
literature trail from Gentry's paper rather than from the
chapter's prose.

\paragraph{The oblivious-computing tradition.}
The trapdoor-paradigm framing of \Cref{sec:fhe-sidebar} sits in a
literature that goes back to early work on secure computation.
\textcite{goldreich-1996}, ``Foundations of Cryptography:
Fragments of a Book'' (already cited in chapter~16 for the
oblivious-Turing-machine model), gives the foundational treatment
of computation under cryptographic constraints, including the
relation between oblivious computation and the secure-multiparty
setting that cipher maps generalize. \textcite{naor-yung-1989},
``Universal One-Way Hash Functions and their Cryptographic
Applications,'' is the older entry point to the trapdoor-function
line of work that the cipher-map abstraction inherits its name
from. The chapter has used ``trapdoor'' in the colloquial sense
rather than the strict cryptographic sense, but the connection to
the formal trapdoor-function literature is real and the
encrypted-search applications of chapter~18 will activate it more
explicitly. The bibliographic-notes treatment is necessarily
brief; the secure-multiparty-computation literature is enormous,
and the reader seeking the full context should consult one of
the standard textbook references in cryptography.


% Exercises
% §17 Exercises
%
% Cross-reference resolution status:
%
% DEFINED here:
%   \label{ex:ch17-1} through \label{ex:ch17-11}
%
% RESOLVED (used here, defined elsewhere):
%   \Cref{def:cipher-map}                     -- §17.1
%   \Cref{sec:bernoulli-error-model}          -- §17.2
%   \Cref{thm:cipher-composition}             -- §17.2
%   \Cref{sec:trapdoor-paradigm}              -- §17.3
%   \Cref{sec:fhe-sidebar}                    -- §17.4
%   \Cref{sec:boolean-trapdoor}               -- §17.5
%   \Cref{def:trapdoor-set}                   -- §17.5
%   \Cref{sec:maximizing-confidentiality}     -- §17.6
%   \Cref{def:ppv}                            -- chapter 7 PPV definition
%   \Cref{def:npv}                            -- chapter 7 NPV definition
%   \Cref{def:perfect-filter-line}            -- chapter 6 §6.4
%   \Cref{sec:cancellation}                   -- chapter 6 §6.3
%   \Cref{thm:composition}                    -- chapter 3 §3.3
%   \Cref{thm:kronecker}                      -- chapter 4 §4.4
%   \Cref{ax:element-indep}                   -- chapter 4 Axiom 1
%   \Cref{sec:bernoulli-data-type}            -- chapter 13 algebraic types
%
% 11 exercises, 3 star / 5 starstar / 3 starstarstar:
% - star x 3: cipher-map PPV/C numerical; composition of channel
%   matrices; FHE vs cipher map in two sentences
% - starstar x 5: multi-query leakage via Kronecker;
%   adaptive attack at O(|A|/(1-fprate)); confidentiality vs
%   space; conjunctive Boolean fprate; cancellation in trapdoor
%   complement-union
% - starstarstar x 3: PPV-based confidentiality bound derivation;
%   trapdoor set as Bernoulli<T> instance; Shannon entropy
%   confidentiality in bits per query

% Difficulty-marker macros (matching chapters 1-16):
\providecommand{\exstar}{$\star$}
\providecommand{\exstarstar}{$\star\star$}
\providecommand{\exstarstarstar}{$\star\star\star$}

\section*{Exercises}
\addcontentsline{toc}{section}{Exercises}

\begin{enumerate}

% -----------------------------------------------------------------------
% WARM-UP  (star x 3): exercises 1, 2, 3
% -----------------------------------------------------------------------

\item \exstar
\label{ex:ch17-1}
\textbf{Cipher-map PPV and confidentiality.}
A trapdoor set is configured with cipher-map rates $\fprate =
0.05$ and $\fnrate = 0$ on a workload with prevalence $\pi =
0.1$.
\begin{enumerate}
  \item[(a)] Substitute into \Cref{def:ppv} to compute $\PPV$.
    Verify your answer numerically.
  \item[(b)] Compute the confidentiality $C = 1 - \PPV$ from
    \Cref{sec:maximizing-confidentiality}. State in one sentence
    what $C$ means operationally: given a server $\top$-report,
    what fraction of the report's information is genuine and
    what fraction protects the plaintext?
\end{enumerate}

\item \exstar
\label{ex:ch17-2}
\textbf{Cipher-map composition.}
Let $Q_1$ and $Q_2$ be the $2 \times 2$ channel matrices of two
cipher maps from \Cref{sec:bernoulli-error-model}, both with
$\fnrate = 0$:
\[
  Q_i \;=\;
  \begin{pmatrix}
    1 - \fprate_i & \fprate_i \\
    0             & 1
  \end{pmatrix},
  \qquad i = 1, 2.
\]
Compute the composed matrix $Q_1 Q_2$ for the chained cipher map.
State the $(\fprate, \fnrate)$ point of the composed map and
verify that $\fnrate = 0$ is preserved under composition.

\item \exstar
\label{ex:ch17-3}
\textbf{FHE versus cipher map: when to use each.}
State in two sentences when a designer should choose FHE
(\Cref{sec:fhe-sidebar}) over a cipher-map approximation and
when the choice should go the other way. Name one application
that motivates each choice. (A two-sentence answer is
sufficient.)

% -----------------------------------------------------------------------
% MEDIUM  (starstar x 5): exercises 4, 5, 6, 7, 8
% -----------------------------------------------------------------------

\item \exstarstar
\label{ex:ch17-4}
\textbf{Multi-query leakage via Kronecker factorization.}
A trapdoor set answers a batch of $m = 10$ queries at
$\fprate = 0.05$ and $\fnrate = 0$, with the queries drawn
independently (chapter~4 Axiom~1 in force; see
\Cref{ax:element-indep}).
\begin{enumerate}
  \item[(a)] Apply the chapter~4 Kronecker theorem
    \Cref{thm:kronecker} to write the joint channel matrix for
    the batch as a Kronecker product of the per-query matrices.
    What is the shape of the joint matrix?
  \item[(b)] Bound the expected number of false-positive
    $\top$-reports across the batch. Compute the numerical
    expectation for $\fprate = 0.05, m = 10$.
  \item[(c)] State in one sentence why iid Kronecker
    factorization does not directly bound the joint
    confidentiality $C$: confidentiality is a posterior, and
    posteriors over multiple queries combine through Bayes'
    rule rather than through matrix products.
\end{enumerate}

\item \exstarstar
\label{ex:ch17-5}
\textbf{Adaptive recovery attack.}
Consider an adversary who queries a trapdoor set holding $|A| = n$
plaintext members at $\fprate = \epsilon$ and $\fnrate = 0$. The
adversary observes a $\top$-report on each query and adapts its
next query based on the report stream.
\begin{enumerate}
  \item[(a)] Construct an adaptive strategy that enumerates the
    plaintext domain in some prescribed order and, on each
    $\top$-report, records the queried element as a candidate
    member.
  \item[(b)] Argue that the expected number of queries needed to
    enumerate $A$ in full is at most $O(n / (1 - \epsilon))$.
    (Hint: each true positive is reported $1 - \fnrate = 1$
    times; each spurious $\top$-report contributes a false
    candidate the adversary discards via a second-pass
    verification.)
  \item[(c)] Discuss in two sentences how the adaptive setting
    relates to the per-query iid analysis: the per-query rates
    are unchanged, but the adversary's posterior over $A$
    grows monotonically with the query count.
\end{enumerate}

\item \exstarstar
\label{ex:ch17-6}
\textbf{Confidentiality versus space.}
Two cipher-map designs for the same plaintext set $A$: design~A
sets $\fprate = 0.01$ and design~B sets $\fprate = 0.001$, both
with $\fnrate = 0$. The workload has $\pi = 0.05$.
\begin{enumerate}
  \item[(a)] Compute $\PPV$ and the confidentiality
    $C = 1 - \PPV$ for both designs from
    \Cref{sec:maximizing-confidentiality}.
  \item[(b)] State which design is more confidential and quantify
    the gap.
  \item[(c)] Using the chapter~5 Bloom-filter space formula,
    state which design uses more storage per inserted element.
    Comment in one sentence on the tension: lower storage is
    the more confidential design.
\end{enumerate}

\item \exstarstar
\label{ex:ch17-7}
\textbf{Conjunctive Boolean leakage.}
A query $x \in A \cap B$ is answered by computing
$\tilde f_A(x) \wedge \tilde f_B(x)$ on two trapdoor sets
(\Cref{sec:boolean-trapdoor}), each at $\fprate = 0.05$ and
$\fnrate = 0$, with independent hash families.
\begin{enumerate}
  \item[(a)] Compute the overall $\fprate$ of the conjunction
    using the chapter~3 composition theorem
    \Cref{thm:composition}.
  \item[(b)] At workload prevalence $\pi = 0.01$ for the
    intersection $A \cap B$, compute the conjunction's $\PPV$
    and confidentiality.
  \item[(c)] Compare with the single-set case at the same $\pi$
    and the same per-set $\fprate$: which is more
    confidential? State in one sentence why.
\end{enumerate}

\item \exstarstar
\label{ex:ch17-8}
\textbf{Cancellation in a trapdoor-set complement-union.}
Apply the chapter~6 cancellation analysis of
\Cref{sec:cancellation} to a trapdoor-set expression of the
form $\complement(\tilde A \cup \tilde B)$, where both
$\tilde A$ and $\tilde B$ have $\fprate = 0.05$ and
$\fnrate = 0$.
\begin{enumerate}
  \item[(a)] By De Morgan, rewrite the expression as
    $\complement \tilde A \cap \complement \tilde B$. State
    the $(\fprate, \fnrate)$ point of each complemented
    operand. (Hint: complementing swaps the two error axes.)
  \item[(b)] Compute the $(\fprate, \fnrate)$ point of the
    intersection of the two complements using the chapter~3
    composition theorem.
  \item[(c)] State whether the resulting cipher map sits on
    the perfect-filter line of \Cref{def:perfect-filter-line}
    or escapes it. (Cancellation lets composed expressions
    leave the line; the question is which axis the resulting
    point sits on, if any.)
\end{enumerate}

% -----------------------------------------------------------------------
% CHALLENGING  (starstarstar x 3): exercises 9, 10, 11
% -----------------------------------------------------------------------

\item \exstarstarstar
\label{ex:ch17-9}
\textbf{Tight lower bound on confidentiality.}
Derive a tight lower bound on the confidentiality
$C(\fprate, \fnrate, \pi)$ of a cipher map, given the
parameters $\fprate, \fnrate, \pi \in [0, 1]$.
\begin{enumerate}
  \item[(a)] Write $C$ as a closed-form rational function of the
    three parameters, starting from the chapter~7 PPV formula
    \Cref{def:ppv} and the identity $C = 1 - \PPV$.
  \item[(b)] At fixed $\pi$, characterize the locus of
    $(\fprate, \fnrate)$ pairs that achieve a target
    confidentiality $C \ge C^{*}$ for some prescribed
    $C^{*} \in (0, 1)$. Show that the locus is a half-plane
    in the Bernoulli square cut by a line through the origin.
  \item[(c)] Specialize to $\fnrate = 0$ (perfect-filter line
    case) and give an explicit lower bound on $\fprate$ as a
    function of $\pi$ and $C^{*}$. Verify your bound on the
    \Cref{sec:maximizing-confidentiality} worked example
    ($\pi = 0.01, C^{*} = 0.832$ requires
    $\fprate \ge \sim 0.05$).
\end{enumerate}

\item \exstarstarstar
\label{ex:ch17-10}
\textbf{Trapdoor set as a $\Bernoulli{T}$ instance.}
Cast a trapdoor set (\Cref{def:trapdoor-set}) as an instance of
the algebraic-types framework of chapter~13.
\begin{enumerate}
  \item[(a)] Identify the underlying plaintext type $T$. For a
    Bloom-filter-backed trapdoor set on integers, what is
    $T$? Be explicit about whether the encryption layer
    changes $T$ or only its representation.
  \item[(b)] Write the $\Bernoulli{T}$ channel matrix that
    governs encrypted membership queries. State which
    Bernoulli-square point the trapdoor set occupies and how
    the encryption layer interacts (or fails to interact) with
    that point.
  \item[(c)] Compose with a second trapdoor set on the same
    type $T$ via the chapter~13 product construction and
    state the $(\fprate, \fnrate)$ point of the composed
    cipher map. Compare with the answer from
    \Cref{thm:cipher-composition}: the two routes should
    agree.
\end{enumerate}

\item \exstarstarstar
\label{ex:ch17-11}
\textbf{Information-theoretic confidentiality in bits per
query.}
Apply chapter~8's Shannon entropy machinery to compute the
information leakage of a cipher map on a Boolean query
``$x \in A$?'' in bits per query, as a function of
$\fprate, \fnrate, \pi$.
\begin{enumerate}
  \item[(a)] Model the adversary's per-query posterior as a
    Bernoulli random variable with parameter $\PPV$ for
    $\top$-reports and parameter $1 - \NPV$ for
    $\bot$-reports. Write the Shannon entropy of the prior
    (a Bernoulli at $\pi$) and the conditional entropies of
    the two posteriors.
  \item[(b)] The per-query leakage in bits is
    $L = H(\pi) - \mathbb{E}_{\text{report}}[H(\text{posterior})]$,
    the prior entropy minus the expected posterior entropy.
    Write $L$ as a closed-form function of
    $\fprate, \fnrate, \pi$.
  \item[(c)] Specialize to the \Cref{sec:maximizing-confidentiality}
    worked example ($\fprate = 0.05, \fnrate = 0, \pi = 0.01$)
    and compute $L$ numerically. Compare with the FHE limit
    ($\fprate = \fnrate = 0$, $L = H(\pi)$, the full
    prior-entropy leakage): the cipher map leaks a small
    fraction of $H(\pi)$ per query, and the framework gives
    a quantitative measure of that fraction.
\end{enumerate}

\end{enumerate}

